\documentclass[aspectratio=169,10pt]{beamer}
\usepackage[brazilian]{babel}

\usetheme{UELdark}

\title{Título da apresentação}
\subtitle{Subtítulo opcional}
\author{Nome do Autor}
\institute{Departamento de Computação\\Universidade Estadual de Londrina}
\date{\today}

% Short forms shown in the footer
\title[Título curto]{Título da apresentação}
\author[Nome do Autor]{Nome do Autor}

\begin{document}

\maketitle

\begin{frame}{Sumário}
  \tableofcontents
\end{frame}

\section{Introdução}

\begin{frame}{Texto e listas}{Um subtítulo de quadro}
  O \alert{texto de destaque} usa o azul do DC em negrito.
  \begin{itemize}
    \item Primeiro item
    \begin{itemize}
      \item Subitem
    \end{itemize}
    \item Segundo item
  \end{itemize}
  \begin{enumerate}
    \item Passo um
    \item Passo dois
  \end{enumerate}
\end{frame}

\section{Blocos e equações}

\begin{frame}{Blocos}
  \begin{block}{Definição}
    Um bloco padrão, com título em verde.
  \end{block}
  \begin{alertblock}{Atenção}
    Um bloco de alerta, em tom neutro.
  \end{alertblock}
  \begin{exampleblock}{Exemplo}
    $\displaystyle \int_0^1 x^2\,dx = \frac{1}{3}$
  \end{exampleblock}
\end{frame}

\begin{frame}[fragile]{Código e tabela}
  \begin{columns}[T]
    \begin{column}{.5\textwidth}
\begin{verbatim}
def f(x):
    return x ** 2
\end{verbatim}
    \end{column}
    \begin{column}{.5\textwidth}
      \centering
      \begin{tabular}{lr}
        \hline
        Item & Valor \\ \hline
        A & 10 \\
        B & 20 \\ \hline
      \end{tabular}
    \end{column}
  \end{columns}
\end{frame}

\obrigado[Perguntas?]

\end{document}
